\subsubsection{Entendendo ponteiros}

Após compreender o conceito de endereços na linguagem C, o entendimento e a 
aplicação de ponteiros se tornam uma questão de prática. Ponteiros são 
variáveis especiais que armazenam o endereço de memória de outras variáveis de 
um tipo definido.

Para esclarecer, vejamos um exemplo prático de como declarar e utilizar ponteiros:

\begin{excode}
    #include <stdio.h>

    int main(void)
    {
        int x = 10;
        int *ptr = &x;
        printf("Valor de ptr: %d (%p)\n", *ptr, ptr);

        return 0;
    }
\end{excode}

Neste exemplo, inicializamos a variável \(x\) com o valor 10. Em seguida, 
declaramos uma variável \texttt{ptr} do tipo \texttt{int *}, indicando que 
\texttt{ptr} é um ponteiro para um inteiro. O símbolo \texttt{*}, posicionado 
entre o tipo e o nome da variável, define \texttt{ptr} como um ponteiro. 
Atribuímos a \texttt{ptr} o endereço de \(x\), utilizando o operador \texttt{\&}.

Para acessar o valor armazenado no endereço apontado por \texttt{ptr}, usamos 
\texttt{*ptr}, onde o operador \texttt{*} é utilizado para 
deferenciar\footnote{Processo de acessar o valor contido dentro de um endereço 
de memória} o ponteiro. A função \texttt{printf} é então utilizada para 
imprimir tanto o valor de \(x\) acessado por meio de \texttt{ptr} quanto o 
endereço de memória armazenado em \texttt{ptr}.

\subsubsection{Acessando variáveis fora de escopo}

Acessar variáveis fora do escopo atual de uma função pode ser útil em muitas 
situações, e os ponteiros facilitam esse processo. Vamos explorar um exemplo 
prático para ilustrar como isso pode ser feito:

\begin{excode}
    #include <stdio.h>

    void foo(int x)
    {
        x = 20;
        printf("foo: valor de x: %d\n", x);
    }

    int main(void)
    {
        int x = 10;
        printf("main: valor de x: %d\n", x);
        foo(x);
        printf("main: valor de x: %d\n", x);

        return 0;
    }
\end{excode}

No código acima, declaramos uma variável \(x\) com o valor 10 na função main, 
imprimimos o valor de \(x\) e, em seguida, chamamos a função \texttt{foo} 
passando \(x\) como argumento. Dentro de \texttt{foo}, definimos o valor de 
\(x\) como 20 e imprimimos seu valor. No entanto, vale observar que a alteração 
feita dentro de \texttt{foo} não afeta o valor de \(x\) na função main, pois 
\texttt{foo} recebe uma cópia do valor de \(x\) e qualquer alteração é feita 
apenas nessa cópia local.

Portanto, ao executar o programa acima é possível notar que o valor de \(x\) 
continua o mesmo na função \texttt{main} mesmo após a execução de \texttt{foo}:

\begin{code}
    main: valor de x: 10
    foo: valor de x: 20
    main: valor de x: 10
\end{code}

Para alterar efetivamente o valor de \(x\) dentro de foo, podemos utilizar 
ponteiros para acessar o endereço de memória da variável \(x\) original. Veja 
como isso pode ser feito:

\begin{excode}
    #include <stdio.h>

    void foo(int *x)
    {
        *x = 20;
        printf("foo: valor de x: %d\n", *x);
    }

    int main(void)
    {
        int x = 10;
        int *ptr = &x;

        printf("main: valor de x: %d\n", x);
        foo(ptr);
        printf("main: valor de x: %d\n", x);

        return 0;
    }
\end{excode}

Neste código, declaramos um ponteiro \texttt{ptr} que guarda o endereço de 
memória de \(x\). Em seguida, passamos esse ponteiro para a função 
\texttt{foo}. Dentro de \texttt{foo}, utilizamos o operador * para acessar o 
valor apontado pelo ponteiro \(x\) e alteramos seu valor para 20. Dessa forma, 
ao executar o programa, o valor de \(x\) também é alterado na função 
\texttt{main}, resultando no seguinte:

\begin{code}
    main: valor de x: 10
    foo: valor de x: 20
    main: valor de x: 20
\end{code}

Esse exemplo demonstra como os ponteiros permitem acessar e modificar variáveis 
fora do escopo atual de uma função, facilitando a manipulação de dados e a 
comunicação entre diferentes partes de um programa.
